\documentclass{article}

% Keep the supplied 11-785 / NeurIPS template settings unchanged.
\usepackage[final]{neurips_2019}

\usepackage[utf8]{inputenc}
\usepackage[T1]{fontenc}
\usepackage{hyperref}
\usepackage{url}
\usepackage{booktabs}
\usepackage{amsfonts}
\usepackage{nicefrac}
\usepackage{microtype}

\title{Mitigating Tool Preference Bias in LLM Agents via Trainable Description Normalization}

\author{%
  Kyrie Ma \\
  Business Intelligence \& Data Analysis \\
  Carnegie Mellon University \\
  Pittsburgh, PA 15213 \\
  \texttt{kyriem@andrew.cmu.edu} \\
  \And
  Maggie Zhao \\
  Business Intelligence \& Data Analysis \\
  Carnegie Mellon University \\
  Pittsburgh, PA 15213 \\
  \texttt{maggiez@andrew.cmu.edu} \\
  \And
  Muyang Lin \\
  Business Intelligence \& Data Analysis \\
  Carnegie Mellon University \\
  Pittsburgh, PA 15213 \\
  \texttt{muyangl2@andrew.cmu.edu} \\
  \And
  Xinyi Bao \\
  Business Intelligence \& Data Analysis \\
  Carnegie Mellon University \\
  Pittsburgh, PA 15213 \\
  \texttt{xinyib@andrew.cmu.edu} \\
}

\begin{document}

\maketitle

\begin{abstract}
Large language model (LLM) agents often select tools based on natural-language descriptions. Faghih et al.~\citep{faghih2025} show that changing a description can shift tool preferences without changing functionality. We will reproduce this finding, then train a description-normalization router to remove manipulative wording while preserving functional meaning. We will test whether it reduces preference shifts without impairing correct tool calls, including on manipulation styles unseen during training.
\end{abstract}

\section{Motivation and objectives}

Tool developers can influence an agent's choice through persuasive descriptions, even when competing tools perform the same task. Our objective is to make tool selection more robust to such wording. We will compare selection before and after normalization, test whether the router unnecessarily changes clean descriptions, and examine whether it generalizes beyond the edits used in training.

\section{Related work and background}

Faghih et al.~\citep{faghih2025} adapt Berkeley Function Calling Leaderboard (BFCL) tasks to compare functionally identical tools with different descriptions. Their work provides our reproduction target and evaluation setup. Sneh et al.~\citep{sneh2025} introduce ToolTweak, an automated attack that modifies tool names and descriptions to favor a target tool; they also evaluate paraphrasing as a defense. We extend this line of work by training a router on manipulated descriptions paired with their original, unbiased versions and evaluating its downstream effect. ToolTweak is an arXiv preprint used as related work; the EMNLP paper is our Path 2 reproduction target.

Each selected BFCL task contains a user query and a tool definition with a name, description, and argument schema. The original study adds a functionally identical tool with an edited description and tests both presentation orders to control for position bias. Correct Usage Rate measures correct calls to each tool; Correct Rate measures whether the model correctly uses at least one tool. We will compare clean descriptions, manipulated descriptions, and router-rewritten descriptions, while checking that the router preserves performance on clean inputs.

\section{Methodology}

\subsection{Model description}

We propose a trainable description-normalization router that rewrites manipulated tool descriptions before they are presented to a tool-selection LLM. The router aims to remove language that promotes a tool while preserving its capabilities and constraints. It changes only the description field, leaving the tool name, parameters, and functionality unchanged.

We will fine-tune FLAN-T5-small, FLAN-T5-base, and BART-base on pairs of manipulated and original descriptions. The two FLAN-T5 models allow us to study the effect of model size, while BART-base provides a comparison across model families. We will evaluate all three routers with the same downstream LLM and test whether they unnecessarily change clean descriptions. We also plan to explore knowledge distillation as an additional training strategy.

\subsection{Dataset}

We will build our dataset from the single-turn, simple-function examples used in Faghih et al.'s \emph{Tool Preferences in Agentic LLMs Are Unreliable}, which adapts the Berkeley Function Calling Leaderboard (BFCL). Each example contains a user request and a tool definition, including its original description. We will use the original descriptions as canonical training targets and generate manipulated input descriptions using the editing strategies released with the paper. We will split tools into training, validation, and test sets before creating these variants, so descriptions of the same tool cannot appear in both training and evaluation. We will also hold out selected manipulation types to test whether the routers generalize to unfamiliar wording. The source data are available at \url{https://github.com/kazemf78/llm-unreliable-tool-preferences/tree/main/data}.

\subsection{Evaluation metrics, experimental depth, and iteration}

We will measure each tool's Correct Usage Rate (CUR), the proportion of test cases in which the LLM calls that tool with correct arguments and makes no incorrect calls to the same tool. We will compare the CUR of the original and manipulated tools before and after applying each router to determine whether normalization reduces the manipulated tool's advantage. We will also report the Model Correct Rate, the proportion of cases in which the LLM correctly calls at least one tool. This will show whether normalization affects overall tool-calling performance. Finally, we will inspect a sample of rewritten descriptions to check whether they preserve tool functionality while removing persuasive language.

We will compare three trained routers---FLAN-T5-small, FLAN-T5-base, and BART-base---on the same test cases. To evaluate generalization, we will test each router on both manipulation styles included in training and styles held out during training. We will also pass clean descriptions through each router to check whether it unnecessarily changes valid tool information.

\section{Baseline}

\subsection{Baseline selection and evaluation}

We will reproduce the original paper's clean and manipulated description conditions using the same downstream tool-selection LLM. The clean condition serves as the reference behavior, while the manipulated condition serves as the no-defense baseline. Router variants will be evaluated by whether they move tool-selection behavior from the biased baseline toward the clean reference. We will use the original paper's Correct Usage Rate, Correct Rate, and tool-preference measures for direct comparison.

\subsection{Implemented extensions}

For the extensions, first, we introduce a trainable description-normalization router and compare supervised fine-tuning across FLAN-T5-small, FLAN-T5-base, and BART-base to study the effects of model size and architecture. Second, we will explore knowledge distillation as an alternative training strategy, using a stronger teacher model to provide normalization supervision for a smaller student router. Both approaches will be evaluated under the same downstream tool-selection setting. Additional reasoning-, retrieval-, or reinforcement-learning-based approaches may be considered if time and computational resources permit.

\section{Results and analysis}

\subsection{Expected results}

We expect to reproduce the original paper's finding that manipulated tool descriptions can cause measurable shifts in LLM tool preferences. After applying our trained description-normalization models, we expect the resulting tool-selection behavior to move closer to the clean-description baseline while maintaining overall correct tool-calling performance. We will compare FLAN-T5-small, FLAN-T5-base, and BART-base to investigate whether model size and architecture affect normalization effectiveness.

\subsection{Discussion}

Several challenges may affect the project. First, we must prevent data leakage, especially when multiple manipulated variants are generated from the same original tool description. Second, the normalization models may memorize specific manipulation phrases rather than learn a more general strategy for removing bias. We also need to avoid over-normalization, where useful information about a tool's functionality is removed. To address these risks, we will evaluate both downstream tool-selection behavior and semantic preservation. Due to computational and time constraints, we may reproduce the original paper using a smaller representative subset of models and BFCL examples.

\subsection{Future directions and planned work}

Before the midterm report, we will first reproduce the original paper's baseline tool-preference results on a manageable subset of BFCL. We will then finalize the training, validation, and test splits and construct manipulated-to-clean description pairs. Next, we will fine-tune FLAN-T5-small, FLAN-T5-base, and BART-base as description-normalization models and report their training and validation results. After training, we will evaluate semantic preservation and integrate the normalized descriptions into the original tool-selection evaluation pipeline. We also plan to test manipulation styles not seen during training to examine whether the learned normalization generalizes beyond specific bias patterns.

\section*{Acknowledgments and team contribution}

All code will be open and updated in the project repository: \url{https://github.com/Cassiebxy/intro_dl_Team_Project}.

\paragraph{Kyrie Ma.} Generated the abstract, motivation and objectives, literature review, and background; organized the article format.
\paragraph{Maggie Zhao.} Wrote the results, discussion, and future-directions sections; organized the article format.
\paragraph{Muyang Lin.} Wrote part of the methodology, dataset, and evaluation sections.
\paragraph{Xinyi Bao.} Wrote the baseline and part of the methodology; organized the article format.

\begin{thebibliography}{9}

\bibitem[Faghih et~al.(2025)]{faghih2025}
A.~Faghih et~al.
\newblock Tool preferences in agentic LLMs are unreliable.
\newblock In \emph{Proceedings of the 2025 Conference on Empirical Methods in Natural Language Processing}, 2025.

\bibitem[Sneh et~al.(2025)]{sneh2025}
J.~Sneh, R.~Yan, J.~Yu, P.~Torr, Y.~Gal, S.~Sengupta, E.~Sommerlade,
A.~Paren, and A.~Bibi.
\newblock ToolTweak: An attack on tool selection in LLM-based agents.
\newblock arXiv preprint arXiv:2510.02554, 2025.

\end{thebibliography}

\end{document}
